\E{dendroschism}{ˌdɛndrəʊˈskɪzəm}{n.}{Gk. δένδρον tree + σχίσμα a split} The forking of a codebase; the division of an open project into two lineages, each claiming to be the trunk. \ex{Every dendroschism is a small evolution, and the fitter branch keeps the stars.} \src{Chron. Fabr.} --- \dv{dendroschismatic} \textit{adj.}; \dv{dendroschist} \textit{n.} one who forks. Cf. \xr{anacalyptism}.

\E{dikaiotachy}{daɪˈkaɪəʊˌtæki}{n.}{Gk. δίκαιος just + τάχος speed} The speed of justice; the rate at which a polity grants permits, approvals and judgments, held by the e/acc to be as much a part of justice as the verdict. \ex{Delayed approval is denied approval: such is the first law of dikaiotachy.} \src{Tract. Cop.} viii.

\E{diaporeusis}{ˌdaɪəpəˈruːsɪs}{n.}{Gk. διαπορεύεσθαι to travel through} Transit between worlds; the crossing of interplanetary space by ships carrying cargo, crews or machines. \sn{1} The voyage itself. \sn{2} \lbl{collect.} The whole commerce of such voyages. --- \dv{diaporeutic} \textit{adj.} Cf. \xr{apoikiast}.

\E{diaskeuomania}{daɪˌæskjuːəʊˈmeɪniə}{n.}{Gk. διασκευή revision, recension + μανία madness} The mania of endless revision; the refusal to ship until the work is perfect, which is to say the refusal to ship. Opp. \xr{ecpempsis}. \ex{The diaskeuomaniac polished the launch until the market had moved to another planet.} \src{Ep. Fabr.} iii. 2. --- \dv{diaskeuomaniac} \textit{n.}

\E{diploergy}{ˈdɪpləʊɜːdʒi}{n.}{Gk. διπλόος double + ἔργον work} The doubling of a civilization's energy supply; the unit by which the anabathmists reckon progress, one diploergy being worth more than ten reforms. \ex{Two diploergies and the debate about scarcity ends of its own accord.} --- \dv{diploergic} \textit{adj.}

\E{dogmatotachy}{ˌdɒɡmətəʊˈtæki}{n.}{Gk. δόγμα opinion, decree + τάχος speed} The speed at which received opinion changes; on the Agora, often faster than the technologies it concerns. \ex{By the dogmatotachy of the Agora, a heresy of Monday is an orthodoxy by Friday and a cliché by Sunday.}

\E{dorocracy}{dəˈrɒkrəsi}{n.}{Gk. δῶρον gift + κράτος rule} \sn{1} The rule of the grant-giver; the government of science by committees of benefaction, whose gifts arrive late and with conditions. \sn{2} \lbl{iron.} Any patronage that enforces caution in the name of support. --- \dv{dorocrat} \textit{n.} Cf. \xr{bouleutophobia}.

\E{drachmogenesis}{ˌdrækməʊˈdʒɛnɪsɪs}{n.}{Gk. δραχμή drachma + γένεσις birth} The first funding of a venture; the seed-round, considered as the moment at which an idea first acquires a metabolism. \ex{At drachmogenesis the founders had a deck, a demo and a disregard for sleep.} \src{Chron. Fabr.} --- \dv{drachmogenetic} \textit{adj.}

\E{dynamotheca}{ˌdaɪnəməʊˈθiːkə}{n., pl. -cae}{Gk. δύναμις power + θήκη case, repository} A battery; a store of power, regarded as the granary of the electric age. \ex{The grid wanted dynamothecae as Egypt wanted granaries.} --- \dv{dynamothecary} \textit{n.} the art of storing energy; \dv{dynamothecal} \textit{adj.}

\E{dyschronia}{dɪsˈkrəʊniə}{n.}{Gk. δυσ- ill + χρόνος time} The mismatch of clocks between institutions and technologies; the condition in which a law drafted in one era is applied to a machine of the next. \ex{Our dyschronia is chronic: the statute walks, the model flies.} \src{Hom. Grad.} xi. --- \dv{dyschronic} \textit{adj.}

\E{ectypomania}{ˌɛktɪpəʊˈmeɪniə}{n.}{Gk. ἔκτυπος worked in relief, copied + μανία madness} The compulsion to copy a successful product; the swarm of imitations that follows every launch. \sn{1} \lbl{derog.} Unoriginal imitation. \sn{2} \lbl{laudatory} The rapid diffusion of a good idea, by which open societies outpace closed ones. --- \dv{ectypomaniac} \textit{n.}

\E{elaterurgy}{ɪˈlætərɜːdʒi}{n.}{Gk. ἐλατήρ driver + -ουργία working} The craft of propulsion; rocketry and its descendants, considered as the art of driving matter out of the gravity-well. \ex{Elaterurgy is the only liturgy that answers prayers upward.} \src{Lib. Ascens.} iv. --- \dv{elaterurge} \textit{n.} a rocket engineer. Cf. \xr{ouranobasis}.

\E{emporotachy}{ˌɛmpərəʊˈtæki}{n.}{Gk. ἐμπόριον market, trading-place + τάχος speed} The velocity of trade; the rate at which goods, contracts and capital change hands, held to be the pulse of a healthy \xr{chrematosphere}.

\E{ephemerotechny}{ɪˌfɛmərəʊˈtɛkni}{n.}{Gk. ἐφήμερος lasting a day + τέχνη art} \sn{1} The art of building things meant to be replaced; the deliberate acceptance that every product is a draft of its successor. \sn{2} \lbl{derog.} Disposable novelty. \ex{The ephemerotechnist ships each version as a gift to the next.} --- \dv{ephemerotechnist} \textit{n.}

\E{epichrysis}{ˌɛpɪˈkraɪsɪs}{n.}{Gk. ἐπί upon + χρυσός gold} A rush of capital toward a fashionable sector; a gold-fever of the venture markets, from which a few great firms and many ruins emerge. \ex{The epichrysis passed; the logisteria it paid for remained.} \src{Chron. Fabr.} --- \dv{epichrystic} \textit{adj.} Cf. \xr{logisterion}.

\E{epicyclist}{ˌɛpɪˈsaɪklɪst}{n.}{Gk. ἐπίκυκλος an epicycle} \lbl{derog.} One who preserves a failing prediction by adding complications to it, as the astronomers preserved the geocentric heavens; said chiefly of an \xr{olethromancer} whose doom has been postponed. --- \dv{epicyclism} \textit{n.}

\E{ergolalia}{ˌɜːɡəʊˈleɪliə}{n.}{Gk. ἔργον work + λαλιά chatter} Talk about work in place of work; the condition of a firm whose meetings concerning a product outnumber its commits. \ex{Their ergolalia filled a calendar and emptied a roadmap.} \src{Ep. Fabr.} iv. 1. --- \dv{ergolalist} \textit{n.} Opp. \xr{ergophile}.

\E{eschatometry}{ˌɛskəˈtɒmɪtri}{n.}{Gk. ἔσχατος last + μέτρον measure} The measurement of the end of the world; the assignment of numerical probabilities to doom, practised at conferences and repeated upon the Agora. \ex{The eschatometry of the incumbents rose with each funding round of their rivals.} \src{Sat. Decel.} x. --- \dv{eschatometer} \textit{n.} one who so measures; \dv{eschatometric} \textit{adj.}

\E{exhypnosis}{ˌɛkshɪpˈnəʊsɪs}{n.}{Gk. ἐξ out + ὕπνος sleep} The awakening of a society from a long stagnation; the moment at which a people remembers that it once built cathedrals and canals. \ex{The reactors were the sign of their exhypnosis.} \src{Hom. Grad.} xii. --- \dv{exhypnotic} \textit{adj.}

\E{exocosmism}{ˌɛksəʊˈkɒzmɪzəm}{n.}{Gk. ἔξω outside + κόσμος world} The orientation of a civilization's ambitions beyond its home world; the conviction that the planet is a cradle, not a house. --- \dv{exocosmist} \textit{n.}; \dv{exocosmic} \textit{adj.} Cf. \xr{siderotropism}.

\E{fenerator}{ˈfɛnəreɪtə}{n.}{L. \la{fenerator} money-lender; Fr. fem. \la{fénératrice}} An investor in risky ventures; one who lends capital against the future. \ex{The fenerator backs ten fools to find one prophet.} \src{Tract. Cop.} ix. --- \dv{feneratrice} \textit{n.} (fem.); \dv{feneration} \textit{n.} venture investing.

\E{fervefaction}{ˌfɜːvɪˈfækʃən}{n.}{L. \la{fervefacere} to make hot} The heating of a sector, market or mind into a state of productive excitement. \sn{1} The early phase of a boom. \sn{2} The arousal of a founder by a new idea at an unreasonable hour. --- \dv{fervefacient} \textit{adj.}

\E{fiduciarch}{fɪˈdjuːʃiɑːk}{n.}{L. \la{fiducia} trust, confidence + Gk. ἀρχός ruler} One whose power rests on the confidence of others; a founder kept in office by belief rather than by revenue. --- \dv{fiduciarchy} \textit{n.} \ex{The fiduciarch must ship before the fiducia runs out.}

\E{fissipary}{ˈfɪsɪpəri}{n.}{L. \la{fissus} split + \la{parere} to bring forth} The spawning of new firms from old ones; the diaspora of engineers who leave a great company to found many small ones. \ex{By fissipary one laboratory became forty.} \src{Chron. Fabr.} --- \dv{fissiparist} \textit{n.} a spin-out founder.

\E{flammigerist}{flæˈmɪdʒərɪst}{n.}{L. \la{flammiger} flame-bearing} One who carries the flame of a project after its founder departs; a maintainer of open code. \ex{Glory to the founders; bread to the flammigerists.} --- \dv{flammigerous} \textit{adj.}

\E{fugacist}{fjuːˈɡeɪsɪst}{n.}{L. \la{fugax} fleeting} One who holds that no competitive advantage endures; that every moat is fugitive and must be dug again each year. \ex{The fugacist sleeps lightly, and builds daily.} --- \dv{fugacism} \textit{n.}

\E{fulcrarch}{ˈfʊlkrɑːk}{n.}{L. \la{fulcrum} bedpost, prop + Gk. ἀρχός ruler} One who rules by leverage; a small team that moves a large world by means of a tool. \ex{Give the fulcrarch a model and a place to stand.} \src{Ep. Fabr.} v. --- \dv{fulcrarchy} \textit{n.}

\E{fabulocracy}{ˌfæbjʊˈlɒkrəsi}{n.}{L. \la{fabula} story + Gk. κράτος rule} Government by narrative; the condition in which the most persuasive story, rather than the most accurate model, determines policy. \ex{In the fabulocracy the doomers had the better plot, the builders the better data.} --- \dv{fabulocrat} \textit{n.} Cf. \xr{memetaxis}.

\E{funambulocracy}{fjuːˌnæmbjʊˈlɒkrəsi}{n.}{L. \la{funambulus} rope-walker + Gk. κράτος rule} Governance as a tightrope-walk between acceleration and restraint; the posture of a regulator who wishes to be praised by both parties. --- \dv{funambulocrat} \textit{n.}

\E{futurescence}{ˌfjuːtʃəˈrɛsəns}{n.}{L. \la{futurus} about to be + -escence, after \textit{adolescence}} The state of becoming the future; the awkward youth of a technology between demonstration and ubiquity. \ex{The machine was in its futurescence, clumsy and immortal.} --- \dv{futurescent} \textit{adj.}

\E{gaiergy}{ˈɡaɪɜːdʒi}{n.}{Gk. Γαῖα earth + ἔργον work} Energy drawn from the planet's interior heat; geothermal power, considered as a gift withheld until the drills were deep enough. --- \dv{gaiergic} \textit{adj.}

\E{geitonophobia}{ˌɡaɪtənəʊˈfəʊbiə}{n.}{Gk. γείτων neighbour + φόβος fear} The fear of new neighbours, whether houses, towers, power-lines or persons; the animating passion of the \xr{oikoclast}. \ex{Geitonophobia built more walls than war.} \src{Sat. Decel.} xi. --- \dv{geitonophobe} \textit{n.}

\E{gelotocracy}{ˌdʒɛləˈtɒkrəsi}{n.}{Gk. γέλως laughter + κράτος rule} Rule by laughter; the power of the meme to make a policy ridiculous before it can be made law. --- \dv{gelotocrat} \textit{n.} Cf. \xr{gelotopoeia}.

\E{gelotopoeia}{ˌdʒɛlətəʊˈpiːə}{n.}{Gk. γελωτοποιία the making of laughter} The craft of meme-making; held upon the Agora to be the most effective form of political argument. \ex{The decels wrote white papers; the accelerationists answered with gelotopoeia, and won.} \src{Agora} --- \dv{gelotopoet} \textit{n.} a maker of memes.

\E{genomorphosis}{ˌdʒiːnəʊmɔːˈfəʊsɪs}{n.}{Gk. γένος race, kind + μόρφωσις a shaping} The editing of the genome; the reshaping of living inheritance by deliberate craft. \ex{After genomorphosis, inheritance became a draft.} --- \dv{genomorphic} \textit{adj.} Cf. \xr{biotachy}.

\E{geophygy}{dʒiˈɒfɪdʒi}{n.}{Gk. γῆ earth + φυγή flight} Flight from the Earth; the departure of the first permanent settlers. \sn{1} The event. \sn{2} \lbl{fig.} Any irreversible step into a frontier. --- \dv{geophygic} \textit{adj.}

\E{glossurgy}{ˈɡlɒsɜːdʒi}{n.}{Gk. γλῶσσα tongue, language + -ουργία working} The craft of instructing machines in natural language; prompt-craft, considered as a branch of rhetoric. \ex{The glossurge speaks, and the machine builds.} --- \dv{glossurge} \textit{n.}

\E{gnosiphagy}{ˌnəʊsɪˈfædʒi}{n.}{Gk. γνῶσις knowledge + -φαγία eating} The devouring of knowledge; the training of a model upon the written record of a civilization. \ex{It ate the libraries, and spoke.} \src{Lib. Ascens.} vi. --- \dv{gnosiphage} \textit{n.}

\E{grammatoclast}{ɡrəˈmætəʊklæst}{n.}{Gk. γράμμα writing, document + κλάστης breaker} A breaker of paperwork; a reformer who abolishes forms. Opp. \xr{grammatocracy}. --- \dv{grammatoclasm} \textit{n.}

\E{gymnotechny}{ˈdʒɪmnəʊtɛkni}{n.}{Gk. γυμνός naked + τέχνη art} The practice of building in public; the exposure of work in progress to the scrutiny of the Agora. \ex{Gymnotechny makes every failure a lesson for strangers.} --- \dv{gymnotechnist} \textit{n.}

\E{halmatism}{ˈhælmətɪzəm}{n.}{Gk. ἅλμα a leap} The leap-frogging of an older technology, as when a nation without wires adopts the wireless at once. --- \dv{halmatic} \textit{adj.} \ex{a halmatic economy, skipping a century in a decade}.

\E{hebdomadist}{hɛbˈdɒmədɪst}{n.}{Gk. ἑβδομάς a week} One who ships weekly; a founder whose calendar is measured in sevens. \ex{The hebdomadist rests on the seventh day only to plan the eighth.} \src{Ep. Fabr.} vi. --- \dv{hebdomadism} \textit{n.}

\E{hecatonergate}{ˌhɛkətəˈnɜːɡeɪt}{n.}{Gk. ἑκατόν a hundred + ἐργάτης worker} An engineer of reputed hundredfold productivity; a figure of legend upon the Agora. \ex{Every hecatonergate is a team of one hundred machines.} \src{Agora}

\E{hetaireiarch}{hɛˈtaɪriɑːk}{n.}{Gk. ἑταιρεία company, club + ἀρχός} The leader of a private circle; the moderator of a group chat in which the future is planned before it is announced. --- \dv{hetaireiarchy} \textit{n.}

\E{hieropausis}{ˌhaɪərəʊˈpɔːsɪs}{n.}{Gk. ἱερός sacred + παῦσις a stopping} A sacred pause; a moratorium invested with the dignity of a sabbath. \ex{They declared a hieropausis, and the rest of the world kept working.} \src{Min. Pausiar.} --- \dv{hieropausal} \textit{adj.} Cf. \xr{pausiarch}.

\E{histurgy}{ˈhɪstɜːdʒi}{n.}{Gk. ἱστός web, tissue + -ουργία working} The engineering of living tissue; the printing of organs. --- \dv{histurge} \textit{n.} Cf. \xr{macrobiotarch}.

\E{hodoplast}{ˈhɒdəʊplæst}{n.}{Gk. ὁδός road + πλαστός formed} A maker of roadmaps; one who draws the path that the engineers will ignore. \ex{The hodoplast draws; the builder deviates.} --- \dv{hodoplasty} \textit{n.}

\E{hoplotechny}{ˈhɒpləʊtɛkni}{n.}{Gk. ὅπλον tool, weapon + τέχνη art} Defensive technology; the craft of armour, shield and deterrent, which accelerates as fast as the threats it answers. --- \dv{hoplotechnist} \textit{n.}

\E{hypsibasis}{hɪpˈsɪbəsɪs}{n.}{Gk. ὑψιβάτης walking on high} The ascent of a rocket; the climb through the atmosphere, watched by millions. \ex{At hypsibasis the Agora held its breath.} --- \dv{hypsibatic} \textit{adj.} Cf. \xr{elaterurgy}.

\E{hydrurgy}{ˈhaɪdrɜːdʒi}{n.}{Gk. ὕδωρ water + -ουργία working} The engineering of water at scale; desalination, irrigation and the greening of deserts by cheap power. \ex{Pleonergy first, then hydrurgy, then gardens where the maps said sand.} \src{Tract. Cop.} x. --- \dv{hydrurge} \textit{n.} Cf. \xr{pleonergy}.
